\subsection{Special Forms: Difference of Cubes} 
Let's see what happens when we multiply $(A-B)(A^2+AB+B^2)$ using the rectangle method.
\begin{lpic}{}{5}
\end{lpic}
Notice that all the \makeblank cancel out! This gives us a special form called a \term{difference of cubes}.
\begin{fact}[Recognizing a Difference of Cubes]
A binomial is a \term{difference of cubes} when:
\begin{itemize*}
\item Each term is a perfect cube monomials (\makeblank[1in] a cubed number, and powers of variables multiples of \makeblanksmall).
\item One is positive, the other negative.
\end{itemize*}
\end{fact}
\begin{Example}Factor each difference of cubes.
\begin{lpic}{}{3}
\regtextleft{0.5}{-1}{$x^3-27={}$};
\end{lpic}
\begin{lpic}{}{3}
\regtextleft{0.5}{-1}{$8x^6y^3-27={}$};
\end{lpic}
\end{Example}
\begin{Note}We can generalize this for \emph{any} powers (not just 2 and 3) but we will only consider 2 and 3 for this class.\end{Note}